% =================================================
% Ученический рабочий лист
% =================================================

\worksheettitle
  {Точки на отрезке}
  {Порядок точек, части отрезка и несколько решений}

\studentnameblock

\begin{checkpointbox}[Входная разминка]
  Точка \(C\) лежит между \(A\) и \(B\). \(AB=11\text{ см}\),
  \(AC=4\text{ см}\). Найдите \(CB\).
  \[
    CB=\answerblank[25mm]-\answerblank[25mm]=\answerblank[32mm].
  \]
\end{checkpointbox}

\section{Порядок точек}

\betweenfigure

\begin{fillbox}[Дополните]
  На рисунке точка \answerblank[13mm] лежит между точками
  \answerblank[13mm] и \answerblank[13mm]. Порядок точек можно записать
  так: \answerblank[38mm].
\end{fillbox}

\begin{practicebox}[Задание 1. Прочитайте записи]
  Для каждой записи назовите точку, которая лежит между двумя другими.
  \begin{enumerate}[label=\alph*),itemsep=0.45em]
    \item \(M-K-N\): между точками \answerblank[13mm] и
      \answerblank[13mm] лежит точка \answerblank[13mm];
    \item \(P-T-S\): между точками \answerblank[13mm] и
      \answerblank[13mm] лежит точка \answerblank[13mm].
  \end{enumerate}
\end{practicebox}

\section{Целое и части}

\wholepartsfigure

\begin{fillbox}[Главное равенство]
  Если \(C\) лежит между \(A\) и \(B\), то
  \[
    AB=\answerblank[20mm]+\answerblank[20mm].
  \]
\end{fillbox}

\begin{practicebox}[Задание 2]
  Точка \(K\) лежит между точками \(M\) и \(N\).
  \begin{enumerate}[label=\alph*),itemsep=0.55em]
    \item \(MK=3\text{ см }8\text{ мм}\), \(KN=5\text{ см }7\text{ мм}\).
      Найдите \(MN\).
    \item \(MN=14\text{ см}\), \(KN=6\text{ см }5\text{ мм}\).
      Найдите \(MK\).
  \end{enumerate}
  \writinglines{2}
\end{practicebox}

\section{Две точки внутри отрезка}

\fourpointsfigure

\begin{fillbox}[Дополните]
  При порядке \(A-C-D-B\):
  \[
    AB=\answerblank[16mm]+\answerblank[16mm]+\answerblank[16mm],
  \]
  поэтому
  \[
    CD=\answerblank[18mm]-\answerblank[18mm]-\answerblank[18mm].
  \]
\end{fillbox}

\begin{practicebox}[Задание 3. Отметьте точки и найдите расстояние]
  \studentcaseonefigure
  \[
    CD=\answerblank[25mm]-\answerblank[25mm]-\answerblank[25mm]
      =\answerblank[30mm].
  \]
\end{practicebox}

\begin{thinkbox}[Задание 4. Как расположены точки?]
  На отрезке \(AB\) отмечены точки \(C\) и \(D\).
  \(AB=15\text{ см}\), \(AC=9\text{ см}\), \(DB=8\text{ см}\).

  Сравните:
  \[
    AC+DB=\answerblank[25mm],\qquad AB=\answerblank[25mm].
  \]
  Какая точка находится левее --- \(C\) или \(D\)? \answerblank[25mm]

  Найдите \(CD\). \answerblank[40mm]
\end{thinkbox}

\section{Расстояния от одного конца}

\sameendpointfigure

\begin{practicebox}[Задание 5]
  Точки \(P\) и \(Q\) лежат на отрезке \(MN\). Известно, что
  \(MP=3\text{ см }5\text{ мм}\), \(MQ=8\text{ см }2\text{ мм}\).

  Какая точка расположена ближе к \(M\)? \answerblank[20mm]

  \[
    PQ=\answerblank[28mm]-\answerblank[28mm]=\answerblank[34mm].
  \]
\end{practicebox}

\section{Два возможных положения}

\twopositionsfigure

\begin{warningbox}[Важно]
  Если известно расстояние от точки \(C\) до точки \(D\), проверьте оба
  направления: влево и вправо от \(C\).
\end{warningbox}

\begin{practicebox}[Задание 6. Рассмотрите два случая]
  На отрезке \(AB\) длиной \(10\text{ см}\) точка \(C\) расположена так,
  что \(AC=6\text{ см}\). Отметьте точку \(D\), если \(CD=2\text{ см}\),
  и найдите \(BD\) во всех возможных случаях.

  \studenttwocasesfigure

  Ответы: \(BD=\answerblank[24mm]\) или \(BD=\answerblank[24mm]\).
\end{practicebox}

\begin{checkpointbox}[Задание 7. Сколько решений?]
  Во всех случаях точка \(D\) должна лежать на отрезке \(AB\),
  \(AB=10\text{ см}\). Укажите количество возможных положений точки \(D\).
  \begin{enumerate}[label=\alph*),itemsep=0.55em]
    \item \(AC=5\text{ см}\), \(CD=3\text{ см}\): \answerblank[18mm];
    \item \(AC=1\text{ см}\), \(CD=3\text{ см}\): \answerblank[18mm];
    \item \(AC=2\text{ см}\), \(CD=11\text{ см}\): \answerblank[18mm].
  \end{enumerate}
\end{checkpointbox}

\begin{reflectionbox}[Итог]
  Перед вычислением нужно установить \answerblank[70mm].
\end{reflectionbox}
